% GENERATED FILE. Edit canonical lessons, then run npm run generate:books.
% canonical-source-hash: fnv1a64:0fea79f78f066ab3
% canonical-lessons: SA-C325-gada, SA-C325-sulam, SA-C325-tunirah, SA-C325-sarathih, SA-C325-vidhyati, SA-R325-review-318-325

\chapter{Weapons of Old}
\label{ch:sa-weapons-of-old}
\hlchaptermodality{325}

\begin{chapteropening}
\textbf{By the end of this chapter:} \emph{I can describe how battles were fought long ago.}

Complete the last lesson of chapter 325: I can describe how battles were fought long ago.
\end{chapteropening}

\section[\texorpdfstring{\sk{गदा} (gadā)}{गदा (gadā)}]{\sk{गदा} (gadā) — a mace; a club}

\label{lesson:SA-C325-gada}



\begin{quote}
Before the new one: say the Sanskrit for calls to mind, then the Sanskrit for tells of.
\end{quote}

\subsection*{You'll want to know: \sk{गदा}}
\textbf{\sk{गदा}} — \emph{gadā} — \textquotedblleft{}a mace; a club\textquotedblright{}.

\begin{quote}
\sk{वीरः} \sk{हस्ते} \sk{गदां} \sk{धारयति।}
\end{quote}

\begin{quote}
\emph{vīraḥ haste gadāṃ dhārayati.}
\end{quote}

\begin{quote}
The hero holds a mace in his hand.
\end{quote}

A feminine noun.

\begin{cousinweb}
\textbf{\sk{गदा}} (\emph{gadā}) MW gives \emph{gadā} as a mace, club or bludgeon, without a derivation, so its origin is uncertain (MW).
\end{cousinweb}

\begin{grammarlens}[title={What you've built}]
One more word for telling what happened.
\end{grammarlens}

\subsection*{Your turn}
\begin{itemize}
\raggedright
  \item \emph{Say it:} \emph{gadā}
  \item \emph{Say it:} \emph{gadā}, once more
  \item \emph{Say it:} \emph{vīraḥ haste gadāṃ dhārayati.} — The hero holds a mace in his hand.
  \item \emph{Recall:} say the Sanskrit for calls to mind, then the Sanskrit for tells of, then say \emph{gadā} again
\end{itemize}

\begin{tcolorbox}[breakable,colback=teal!4,colframe=teal!35!black,title={Before you move on}]
What does \sk{गदा} mean? (\textquotedblleft{}A mace; a club\textquotedblright{}.) Say it once more.
\end{tcolorbox}
\section[\texorpdfstring{\sk{शूलम्} (śūlam)}{शूलम् (śūlam)}]{\sk{शूलम्} (śūlam) — a spear; a trident}

\label{lesson:SA-C325-sulam}



\begin{quote}
Before the new one: say the Sanskrit for tells of, then the Sanskrit for a mace.
\end{quote}

\subsection*{You'll want to know: \sk{शूलम्}}
\textbf{\sk{शूलम्}} — \emph{śūlam} — \textquotedblleft{}a spear; a trident\textquotedblright{}.

\begin{quote}
\sk{सैनिकः} \sk{शूलेन} \sk{युध्यते।}
\end{quote}

\begin{quote}
\emph{sainikaḥ śūlena yudhyate.}
\end{quote}

\begin{quote}
The soldier fights with a spear.
\end{quote}

A neuter noun.

\begin{cousinweb}
\textbf{\sk{शूलम्}} (\emph{śūlam}) MW gives \emph{śūla} as a sharp iron pin or stake, and any pointed dart, lance, pike or spear, especially the trident of Śiva, without a derivation, so its origin is uncertain (MW).
\end{cousinweb}

\begin{grammarlens}[title={What you've built}]
One more word for telling what happened.
\end{grammarlens}

\subsection*{Your turn}
\begin{itemize}
\raggedright
  \item \emph{Say it:} \emph{śūlam}
  \item \emph{Say it:} \emph{śūlam}, once more
  \item \emph{Say it:} \emph{sainikaḥ śūlena yudhyate.} — The soldier fights with a spear.
  \item \emph{Recall:} say the Sanskrit for tells of, then the Sanskrit for a mace, then say \emph{śūlam} again
\end{itemize}

\begin{tcolorbox}[breakable,colback=teal!4,colframe=teal!35!black,title={Before you move on}]
What does \sk{शूलम्} mean? (\textquotedblleft{}A spear; a trident\textquotedblright{}.) Say it once more.
\end{tcolorbox}
\section[\texorpdfstring{\sk{तूणीरः} (tūṇīraḥ)}{तूणीरः (tūṇīraḥ)}]{\sk{तूणीरः} (tūṇīraḥ) — a quiver}

\label{lesson:SA-C325-tunirah}



\begin{quote}
Before the new one: say the Sanskrit for a mace, then the Sanskrit for a spear.
\end{quote}

\subsection*{You'll want to know: \sk{तूणीरः}}
\textbf{\sk{तूणीरः}} — \emph{tūṇīraḥ} — \textquotedblleft{}a quiver\textquotedblright{}.

\begin{quote}
\sk{वीरस्य} \sk{तूणीरे} \sk{बाणाः} \sk{सन्ति।}
\end{quote}

\begin{quote}
\emph{vīrasya tūṇīre bāṇāḥ santi.}
\end{quote}

\begin{quote}
There are arrows in the hero's quiver.
\end{quote}

A masculine noun.

\begin{cousinweb}
\textbf{\sk{तूणीरः}} (\emph{tūṇīraḥ}) MW gives \emph{tūṇīra} as the same as \emph{tūṇa}, a quiver, without a further derivation, so its origin is uncertain (MW).
\end{cousinweb}

\begin{grammarlens}[title={What you've built}]
One more word for telling what happened.
\end{grammarlens}

\subsection*{Your turn}
\begin{itemize}
\raggedright
  \item \emph{Say it:} \emph{tūṇīraḥ}
  \item \emph{Say it:} \emph{tūṇīraḥ}, once more
  \item \emph{Say it:} \emph{vīrasya tūṇīre bāṇāḥ santi.} — There are arrows in the hero's quiver.
  \item \emph{Recall:} say the Sanskrit for a mace, then the Sanskrit for a spear, then say \emph{tūṇīraḥ} again
\end{itemize}

\begin{tcolorbox}[breakable,colback=teal!4,colframe=teal!35!black,title={Before you move on}]
What does \sk{तूणीरः} mean? (\textquotedblleft{}A quiver\textquotedblright{}.) Say it once more.
\end{tcolorbox}
\section[\texorpdfstring{\sk{सारथिः} (sārathiḥ)}{सारथिः (sārathiḥ)}]{\sk{सारथिः} (sārathiḥ) — a charioteer}

\label{lesson:SA-C325-sarathih}



\begin{quote}
Before the new one: say the Sanskrit for a spear, then the Sanskrit for a quiver.
\end{quote}

\subsection*{You'll want to know: \sk{सारथिः}}
\textbf{\sk{सारथिः}} — \emph{sārathiḥ} — \textquotedblleft{}a charioteer\textquotedblright{}.

\begin{quote}
\sk{सारथिः} \sk{रथं} \sk{वेगेन} \sk{चालयति।}
\end{quote}

\begin{quote}
\emph{sārathiḥ rathaṃ vegena cālayati.}
\end{quote}

\begin{quote}
The charioteer drives the chariot fast.
\end{quote}

A masculine noun.

\begin{cousinweb}
\textbf{\sk{सारथिः}} (\emph{sārathiḥ}) is built on \emph{sa-ratha}, \textquotedblleft{}with a chariot\textquotedblright{}, from \emph{ratha}, \textquotedblleft{}a chariot\textquotedblright{}; MW gives \emph{sārathi} as a charioteer, a driver (MW).
\end{cousinweb}

\begin{grammarlens}[title={What you've built}]
One more word for telling what happened.
\end{grammarlens}

\subsection*{Your turn}
\begin{itemize}
\raggedright
  \item \emph{Say it:} \emph{sārathiḥ}
  \item \emph{Say it:} \emph{sārathiḥ}, once more
  \item \emph{Say it:} \emph{sārathiḥ rathaṃ vegena cālayati.} — The charioteer drives the chariot fast.
  \item \emph{Recall:} say the Sanskrit for a spear, then the Sanskrit for a quiver, then say \emph{sārathiḥ} again
\end{itemize}

\begin{tcolorbox}[breakable,colback=teal!4,colframe=teal!35!black,title={Before you move on}]
What does \sk{सारथिः} mean? (\textquotedblleft{}A charioteer\textquotedblright{}.) Say it once more.
\end{tcolorbox}
\section[\texorpdfstring{\sk{विध्यति} (vidhyati)}{विध्यति (vidhyati)}]{\sk{विध्यति} (vidhyati) — pierces}

\label{lesson:SA-C325-vidhyati}



\begin{quote}
Before the new one: say the Sanskrit for a quiver, then the Sanskrit for a charioteer.
\end{quote}

\subsection*{You'll want to know: \sk{विध्यति}}
\textbf{\sk{विध्यति}} — \emph{vidhyati} — \textquotedblleft{}pierces\textquotedblright{}.

\begin{quote}
\sk{बाणः} \sk{लक्ष्यं} \sk{विध्यति।}
\end{quote}

\begin{quote}
\emph{bāṇaḥ lakṣyaṃ vidhyati.}
\end{quote}

\begin{quote}
The arrow pierces the target.
\end{quote}

Past (he …): \sk{सः} \sk{लक्ष्यं} \sk{विद्धवान्}, \emph{saḥ lakṣyaṃ viddhavān}.

\begin{cousinweb}
\textbf{\sk{विध्यति}} (\emph{vidhyati}) comes from the root √vyadh, \textquotedblleft{}to pierce, wound\textquotedblright{}; MW gives \emph{viddha} as its past participle, pierced, stabbed, wounded (MW).
\end{cousinweb}

\begin{grammarlens}[title={What you've built}]
One more word for telling what happened.
\end{grammarlens}

\subsection*{Your turn}
\begin{itemize}
\raggedright
  \item \emph{Say it:} \emph{vidhyati}
  \item \emph{Say it:} \emph{vidhyati}, once more
  \item \emph{Say it:} \emph{bāṇaḥ lakṣyaṃ vidhyati.} — The arrow pierces the target.
  \item \emph{Recall:} say the Sanskrit for a quiver, then the Sanskrit for a charioteer, then say \emph{vidhyati} again
\end{itemize}

\begin{tcolorbox}[breakable,colback=teal!4,colframe=teal!35!black,title={Before you move on}]
What does \sk{विध्यति} mean? (\textquotedblleft{}Pierces\textquotedblright{}.) Say it once more.
\end{tcolorbox}
\section[(dialogue)]{Looking back: from The Army Marches to Weapons of Old}

\label{lesson:SA-R325-review-318-325}



\begin{quote}
Say the Sanskrit for an army and a warrior.
\end{quote}

\subsection*{Your turn}
Say these aloud:

\begin{itemize}
\raggedright
  \item an army, a warrior, a war, sets out, fights
  \item a victory, a defeat, a captive, retreats, overpowers
  \item the subjects, a slave, a revolt, drives out, freedom
  \item a coronation, a family priest, a sacred verse, anoints, pronounces
  \item an ancestor, a family line, a clan, founds, introduces
  \item a stupa, a rampart, the past, digs up, brings to light
  \item long ago, a chronicle, an ancient tale, calls to mind, tells of
  \item a mace, a spear, a quiver, a charioteer, pierces
\end{itemize}

\begin{tcolorbox}[breakable,colback=teal!4,colframe=teal!35!black,title={Before you move on}]
An army? (\textbf{\sk{सेना}}, \emph{senā}.) A warrior? (\textbf{\sk{योधः}}, \emph{yodhaḥ}.) A war? (\textbf{\sk{युद्धम्}}, \emph{yuddham}.) Sets out? (\textbf{\sk{प्रयाति}}, \emph{prayāti}.) Fights? (\textbf{\sk{युध्यते}}, \emph{yudhyate}.) A victory? (\textbf{\sk{विजयः}}, \emph{vijayaḥ}.) A defeat? (\textbf{\sk{पराजयः}}, \emph{parājayaḥ}.) A captive? (\textbf{\sk{बन्दी}}, \emph{bandī}.) Retreats? (\textbf{\sk{अपसरति}}, \emph{apasarati}.) Overpowers? (\textbf{\sk{अभिभवति}}, \emph{abhibhavati}.) The subjects? (\textbf{\sk{प्रजा}}, \emph{prajā}.) A slave? (\textbf{\sk{दासः}}, \emph{dāsaḥ}.) A revolt? (\textbf{\sk{विप्लवः}}, \emph{viplavaḥ}.) Drives out? (\textbf{\sk{निष्कासयति}}, \emph{niṣkāsayati}.) Freedom? (\textbf{\sk{स्वातन्त्र्यम्}}, \emph{svātantryam}.) A coronation? (\textbf{\sk{अभिषेकः}}, \emph{abhiṣekaḥ}.) A family priest? (\textbf{\sk{पुरोहितः}}, \emph{purohitaḥ}.) A sacred verse? (\textbf{\sk{मन्त्रः}}, \emph{mantraḥ}.) Anoints? (\textbf{\sk{अभिषिञ्चति}}, \emph{abhiṣiñcati}.) Pronounces? (\textbf{\sk{उच्चारयति}}, \emph{uccārayati}.) An ancestor? (\textbf{\sk{पूर्वजः}}, \emph{pūrvajaḥ}.) A family line? (\textbf{\sk{वंशः}}, \emph{vaṃśaḥ}.) A clan? (\textbf{\sk{कुलम्}}, \emph{kulam}.) Founds? (\textbf{\sk{संस्थापयति}}, \emph{saṃsthāpayati}.) Introduces? (\textbf{\sk{प्रवर्तयति}}, \emph{pravartayati}.) A stupa? (\textbf{\sk{स्तूपः}}, \emph{stūpaḥ}.) A rampart? (\textbf{\sk{प्राचीरम्}}, \emph{prācīram}.) The past? (\textbf{\sk{अतीतम्}}, \emph{atītam}.) Digs up? (\textbf{\sk{उत्खनति}}, \emph{utkhanati}.) Brings to light? (\textbf{\sk{आविष्करोति}}, \emph{āviṣkaroti}.) Long ago? (\textbf{\sk{पुरा}}, \emph{purā}.) A chronicle? (\textbf{\sk{इतिवृत्तम्}}, \emph{itivṛttam}.) An ancient tale? (\textbf{\sk{पुराणम्}}, \emph{purāṇam}.) Calls to mind? (\textbf{\sk{अनुस्मरति}}, \emph{anusmarati}.) Tells of? (\textbf{\sk{कीर्तयति}}, \emph{kīrtayati}.) A mace? (\textbf{\sk{गदा}}, \emph{gadā}.) A spear? (\textbf{\sk{शूलम्}}, \emph{śūlam}.) A quiver? (\textbf{\sk{तूणीरः}}, \emph{tūṇīraḥ}.) A charioteer? (\textbf{\sk{सारथिः}}, \emph{sārathiḥ}.) Pierces? (\textbf{\sk{विध्यति}}, \emph{vidhyati}.)
\end{tcolorbox}
